% ==========================================================
% titles/title-09.tex -- Title IX: Internal Security and Justice
% ==========================================================

\ConstTitle{Internal Security and Justice}{title:internal-security}

\Article{Internal Security Principles}{art:internal-security}

\ArtSection{No Secret Police}{art:internal-security:s-no-secret-police}
All security operations are subject to real-time oversight by rotating \gls{gls:citizen-oversight-panels}s (Article \ref{art:citizen-oversight-panels}). The existence of all investigations must be publicly logged in real time through the \gls{gls:transparency-engine}, though content may be withheld during active operations.

\ArtSection{No Mass Surveillance}{art:internal-security:s-no-mass-surveillance}
Security bodies are authorized to investigate specific credible threats, not to conduct surveillance of the general population. Mass collection of communications, location data, or behavioral data on residents is constitutionally prohibited and may not be enabled by any emergency declaration.

\ArtSection{Separation of Investigation and Enforcement}{art:internal-security:s-separation-investigation-enforcement}
The body that investigates potential threats may not be the same body that acts on them. Investigation findings are submitted to the \gls{gls:security-enforcement-review-panel}, defined in \ref{art:internal-security:s-security-enforcement-review}, which independently authorizes specific enforcement actions. No single body holds both the power to identify threats and the power to act against them.

\ArtSection{The Investigation Body}{art:internal-security:s-investigation-body}
The \gls{gls:security-investigation-governing-board} is a standard body of nine members serving three-year terms, selected from the \gls{gls:qualification-commons} pool for security investigation competency. The Board sets investigation priorities, authorizes opening and closing of investigations by majority vote, and reviews all findings before submission to the \gls{gls:security-enforcement-review-panel}. A professional investigative staff implements authorized investigations under the direction of a Staff Director selected through the \gls{gls:qualification-commons}, serving five-year terms renewable subject to the reconfirmation mechanism of Article \ref{art:governance-principles}. The Staff Director has operational independence within Board-authorized investigations; any Board attempt to direct specific investigative tactics rather than priorities is automatically logged through the \gls{gls:transparency-engine} as a potential constitutional violation. Content of active investigations may be withheld from the \gls{gls:transparency-engine} by Board majority vote for up to one year, renewable once by supermajority. The \gls{gls:citizen-oversight-panels} (Article \ref{art:citizen-oversight-panels}) embedded in this body has full access to all content including withheld material throughout.

\ArtSection{Protection of Governance Participants}{art:internal-security:s-protection-governance-participants}
Security services may not investigate legislators, judges, council members, or other governance participants without authorization from the \gls{gls:constitutional-court} with full transparency to the relevant \gls{gls:citizen-oversight-panels}. 

\ArtSection{The Security Enforcement Review Panel}{art:internal-security:s-security-enforcement-review}
The \gls{gls:security-enforcement-review-panel} is a standard body of nine members serving three-year terms, with no other governance role. The Panel receives investigation findings and independently authorizes or declines specific enforcement actions. Authorization requires five of nine members. The Panel may not conduct investigations itself. All authorization decisions are logged in real time by the \gls{gls:transparency-engine}. 

\ArtSection{Community-Based Enforcement}{art:internal-security:s-community-based-enforcement}
Enforcement of laws at the local level is community-based, with officers selected from and accountable to the communities they serve, subject to immediate community recall. Recall is achieved via simple majority vote through the \gls{gls:transparency-engine}.

\Article{Justice and Restoration}{art:justice-restoration}

\ArtSection{Restorative Principle}{art:justice-restoration:s-restorative-principle}
The Commonwealth's justice system is oriented toward restoration of harm, community accountability, and addressing the conditions that produce harmful behavior. Punishment as an end in itself, and incarceration as a default response to social problems, are inconsistent with the anti-domination principle.

\ArtSection{Abolition of the Carceral Default}{art:justice-restoration:s-abolition-carceral-default}
Mass incarceration is a form of domination. Incarceration is reserved for genuine public safety necessity or demand from the surviving party, and is subject to strict time limits, humane conditions, and regular review by independent \gls{gls:citizen-oversight-panels}s (Article \ref{art:citizen-oversight-panels}).

\ArtSection{Community Justice Bodies}{art:justice-restoration:s-community-justice-bodies}
Restorative justice is implemented through \gls{gls:community-justice-bodies} embedded within local democratic governance under Article \ref{art:local-democracy}. Members are selected by local lottery from qualified pools on rotating terms not to exceed two years. Terms are staggered so that no more than half a Body's membership rotates out in any single selection cycle, and no person may serve consecutive terms without an intervening term served by another resident. \gls{gls:community-justice-bodies} facilitate restorative processes, connect parties to support resources, and implement agreed accountability measures within the minimum standards established by the \gls{gls:justice-standards-council} (\ref{art:justice-restoration:s-justice-standards-council}). Members may not have personal relationships to either party that would compromise impartiality, assessed by the local \gls{gls:citizen-oversight-panels}.

\ArtSection{The Justice Standards Council}{art:justice-restoration:s-justice-standards-council}
A \gls{gls:justice-standards-council}, constituted as an \gls{gls:implementation-councils} under Article \ref{art:implementation-councils}, sets and maintains minimum standards for restorative justice processes Commonwealth-wide. It does not adjudicate individual cases. Its mandate is standard-setting, monitoring, training resource provision, and annual public assessment of restorative justice implementation across communities. All standards are published in full through the \gls{gls:transparency-engine}.

\ArtSection{Minimum Standards}{art:justice-restoration:s-minimum-standards}
The following minimum standards apply to all restorative justice processes without exception and may be raised but not lowered by the \gls{gls:justice-standards-council}:
\begin{itemize}
  \item No restorative process may proceed without the genuine informed consent of the person who experienced harm. Consent may be withdrawn at any point without penalty.
  \item Survivors retain the absolute right to exit a restorative process at any point and to request a carceral response instead, without this choice being used against them in any subsequent proceeding.
  \item Persons who have caused harm may not be coerced into restorative processes as a condition of avoiding harsher consequences in ways that undermine the genuine voluntariness of their participation.
\end{itemize}

All restorative processes are logged by existence through the \gls{gls:transparency-engine}. Process content is protected to preserve integrity and participant privacy. Outcomes and compliance with minimum standards are public.

\ArtSection{Escalation Mechanism}{art:justice-restoration:s-escalation-mechanism}
A survivor who believes a \gls{gls:community-justice-bodies} process has been compromised by local power dynamics may escalate to a regional \gls{gls:justice-standards-council} (\ref{art:justice-restoration:s-justice-standards-council}) review. This review assesses whether minimum standards were met, not the substantive outcome. If minimum standards were not met, the process is void. The survivor may then choose a fresh community process or a carceral response.

\ArtSection{Intimate Partner and Domestic Violence}{art:justice-restoration:s-intimate-partner-domestic}
Restorative justice mechanisms for intimate partner and domestic violence prioritize the agency of the survivor above all other considerations. Survivors determine the form of accountability process. Carceral responses are available whenever survivors choose them. No community norm, facilitation preference, or \gls{gls:justice-standards-council} guidance may be used to discourage survivors from choosing carceral responses.

\ArtSection{Secure Facilities}{art:justice-restoration:s-secure-facilities}
Where incarceration is necessary, secure facilities are governed by the \gls{gls:justice-standards-council} with mandatory independent oversight by rotating \gls{gls:citizen-oversight-panels}s (Article \ref{art:citizen-oversight-panels}). Panels have full and unrestricted access to all facility conditions, records, and residents. Facility conditions must meet the Material Sufficiency Floor without exception, consistent with Article \ref{art:matsuff}, \ref{art:matsuff:s-universality}. All incarceration is subject to strict time limits and mandatory periodic review by independent citizen panels.

\ArtSection{Voluntary Labor by Incarcerated Persons}{art:justice-restoration:s-voluntary-labor-incarcerated}
Incarcerated persons may voluntarily perform wage labor within secure facilities. Participation is strictly voluntary and may not be made a condition of any benefit, privilege, reduced term, or favorable treatment within the facility. Refusal to participate may not result in any penalty, reduction of rights, or adverse treatment of any kind.

Labor performed by incarcerated persons serves only Commons institutions or the person's own skill development. No cooperative, private entity, or commercial operation may benefit from the labor of incarcerated persons.

Wages are set by the \gls{gls:justice-standards-council} (\ref{art:justice-restoration:s-justice-standards-council}) at no less than the essential sector compensation floor established in Article \ref{art:essential-labor}, \ref{art:essential-labor:s-essential-labor}. Incarcerated persons determine whether wages are held as a reintegration fund payable upon release, directed to dependents during incarceration, or divided between both purposes. The \gls{gls:justice-standards-council} may not alter the wage floor downward; it may raise it.

All labor arrangements are logged through the \gls{gls:transparency-engine}. The \gls{gls:citizen-oversight-panels} (Article \ref{art:citizen-oversight-panels}) embedded in the \gls{gls:justice-standards-council} reviews labor conditions annually and publishes findings publicly.